% GENERATED FILE. Edit canonical lessons, then run npm run generate:books.
% canonical-source-hash: fnv1a64:6c194a41ff026d42
% canonical-lessons: LA-C515-vultus, LA-C515-facies, LA-C515-crinis, LA-C515-supercilium, LA-C515-palpebra

\chapter{The face}
\label{ch:la-the-face}
\hlchaptermodality{515}

\begin{chapteropening}
\textbf{By the end of this chapter:} \emph{I can describe a face.}

Complete the last lesson of chapter 515: I can describe a face.
\end{chapteropening}

\section[vultus, vultūs]{vultus, vultūs — face, expression}

\label{lesson:LA-C515-vultus}



\begin{quote}
Before the new one: say the Latin for indignation, then the Latin for to displease.
\end{quote}

\subsection*{You'll want to know: vultus, vultūs}
\textbf{vultus, vultūs} — \textquotedblleft{}face, expression\textquotedblright{}.

\begin{quote}
Vultus mātris laetus est.
\end{quote}

\begin{quote}
\emph{Mother's face is happy.}
\end{quote}

Masculine, fourth declension. The face as it shows feeling and mood.

\begin{cousinweb}
\textbf{vultus} is a native Latin noun; EtymDB links it with a Proto-Indo-European root meaning \textquotedblleft{}to see\textquotedblright{}, while Lewis \& Short calls it doubtful, so the origin is not settled.
\end{cousinweb}

\begin{grammarlens}[title={What you've built}]
One more word for telling, explaining and describing at B1.
\end{grammarlens}

\subsection*{Your turn}
\begin{itemize}
\raggedright
  \item \emph{Say it:} \emph{vultus, vultūs}
  \item \emph{Say it:} \emph{vultus, vultūs}, once more
  \item \emph{Say it:} a sentence of your own with \emph{vultus, vultūs}
  \item \emph{Recall:} say the Latin for indignation, then the Latin for to displease, then say \emph{vultus, vultūs} again
\end{itemize}

\begin{tcolorbox}[breakable,colback=teal!4,colframe=teal!35!black,title={Before you move on}]
What does vultus, vultūs mean? (\textquotedblleft{}Face, expression\textquotedblright{}.) Say it once more.
\end{tcolorbox}
\section[faciēs, faciēī]{faciēs, faciēī — appearance, shape; face}

\label{lesson:LA-C515-facies}



\begin{quote}
Before the new one: say the Latin for to displease, then the Latin for face.
\end{quote}

\subsection*{You'll want to know: faciēs, faciēī}
\textbf{faciēs, faciēī} — \textquotedblleft{}appearance, shape; face\textquotedblright{}.

\begin{quote}
Puer faciem in aquā lavat.
\end{quote}

\begin{quote}
\emph{The boy washes his face in the water.}
\end{quote}

Feminine, fifth declension. The shape and look of a face, or of anything; \emph{vultus} is the face as it shows feeling. English \emph{face}.

\begin{cousinweb}
\textbf{faciēs} is usually traced to the root of \emph{faciō} \textquotedblleft{}to make\textquotedblright{} (EtymDB): the make of a thing, its shape; Lewis \& Short links it with \emph{fax} \textquotedblleft{}torch\textquotedblright{} instead. English \emph{face} and \emph{surface} come from it.
\end{cousinweb}

\begin{grammarlens}[title={What you've built}]
One more word for telling, explaining and describing at B1.
\end{grammarlens}

\subsection*{Your turn}
\begin{itemize}
\raggedright
  \item \emph{Say it:} \emph{faciēs, faciēī}
  \item \emph{Say it:} \emph{faciēs, faciēī}, once more
  \item \emph{Say it:} a sentence of your own with \emph{faciēs, faciēī}
  \item \emph{Recall:} say the Latin for to displease, then the Latin for face, then say \emph{faciēs, faciēī} again
\end{itemize}

\begin{tcolorbox}[breakable,colback=teal!4,colframe=teal!35!black,title={Before you move on}]
What does faciēs, faciēī mean? (\textquotedblleft{}Appearance, shape; face\textquotedblright{}.) Say it once more.
\end{tcolorbox}
\section[crīnis, crīnis]{crīnis, crīnis — hair, a lock of hair}

\label{lesson:LA-C515-crinis}



\begin{quote}
Before the new one: say the Latin for face, then the Latin for appearance.
\end{quote}

\subsection*{You'll want to know: crīnis, crīnis}
\textbf{crīnis, crīnis} — \textquotedblleft{}hair, a lock of hair\textquotedblright{}.

\begin{quote}
Crīnēs puellae longī et flāvī sunt.
\end{quote}

\begin{quote}
\emph{The girl's hair is long and golden.}
\end{quote}

Masculine, third declension; usually plural. A word of poetry and fine prose; \emph{capillus} is the everyday word.

\begin{cousinweb}
\textbf{crīnis} is a native Latin noun; Lewis \& Short and EtymDB compare \emph{crista} \textquotedblleft{}crest\textquotedblright{}, but the origin is not settled.
\end{cousinweb}

\begin{grammarlens}[title={What you've built}]
One more word for telling, explaining and describing at B1.
\end{grammarlens}

\subsection*{Your turn}
\begin{itemize}
\raggedright
  \item \emph{Say it:} \emph{crīnis, crīnis}
  \item \emph{Say it:} \emph{crīnis, crīnis}, once more
  \item \emph{Say it:} a sentence of your own with \emph{crīnis, crīnis}
  \item \emph{Recall:} say the Latin for face, then the Latin for appearance, then say \emph{crīnis, crīnis} again
\end{itemize}

\begin{tcolorbox}[breakable,colback=teal!4,colframe=teal!35!black,title={Before you move on}]
What does crīnis, crīnis mean? (\textquotedblleft{}Hair, a lock of hair\textquotedblright{}.) Say it once more.
\end{tcolorbox}
\section[supercilium, superciliī]{supercilium, superciliī — eyebrow}

\label{lesson:LA-C515-supercilium}



\begin{quote}
Before the new one: say the Latin for appearance, then the Latin for hair.
\end{quote}

\subsection*{You'll want to know: supercilium, superciliī}
\textbf{supercilium, superciliī} — \textquotedblleft{}eyebrow\textquotedblright{}.

\begin{quote}
Senex supercilia alba habet.
\end{quote}

\begin{quote}
\emph{The old man has white eyebrows.}
\end{quote}

Neuter, second declension. Also a frown, and so pride: \emph{supercilium} — haughtiness. English \emph{supercilious}.

\begin{cousinweb}
\textbf{supercilium} is built on \emph{super} \textquotedblleft{}above\textquotedblright{} and \emph{cilium} \textquotedblleft{}eyelid\textquotedblright{} (EtymDB; Lewis \& Short). English \emph{supercilious} comes from it, from a raised eyebrow.
\end{cousinweb}

\begin{grammarlens}[title={What you've built}]
One more word for telling, explaining and describing at B1.
\end{grammarlens}

\subsection*{Your turn}
\begin{itemize}
\raggedright
  \item \emph{Say it:} \emph{supercilium, superciliī}
  \item \emph{Say it:} \emph{supercilium, superciliī}, once more
  \item \emph{Say it:} a sentence of your own with \emph{supercilium, superciliī}
  \item \emph{Recall:} say the Latin for appearance, then the Latin for hair, then say \emph{supercilium, superciliī} again
\end{itemize}

\begin{tcolorbox}[breakable,colback=teal!4,colframe=teal!35!black,title={Before you move on}]
What does supercilium, superciliī mean? (\textquotedblleft{}Eyebrow\textquotedblright{}.) Say it once more.
\end{tcolorbox}
\section[palpebra, palpebrae]{palpebra, palpebrae — eyelid}

\label{lesson:LA-C515-palpebra}



\begin{quote}
Before the new one: say the Latin for hair, then the Latin for eyebrow.
\end{quote}

\subsection*{You'll want to know: palpebra, palpebrae}
\textbf{palpebra, palpebrae} — \textquotedblleft{}eyelid\textquotedblright{}.

\begin{quote}
Puer fessus palpebrās claudit.
\end{quote}

\begin{quote}
\emph{The tired boy closes his eyelids.}
\end{quote}

Feminine, first declension; usually plural. English \emph{palpebral}, in medicine.

\begin{cousinweb}
\textbf{palpebra} is built on \emph{palpō} \textquotedblleft{}to stroke, touch softly\textquotedblright{} (EtymDB), with an ending for a tool: the soft lid that strokes the eye.
\end{cousinweb}

\begin{grammarlens}[title={What you've built}]
One more word for telling, explaining and describing at B1.
\end{grammarlens}

\subsection*{Your turn}
\begin{itemize}
\raggedright
  \item \emph{Say it:} \emph{palpebra, palpebrae}
  \item \emph{Say it:} \emph{palpebra, palpebrae}, once more
  \item \emph{Say it:} a sentence of your own with \emph{palpebra, palpebrae}
  \item \emph{Recall:} say the Latin for hair, then the Latin for eyebrow, then say \emph{palpebra, palpebrae} again
\end{itemize}

\begin{tcolorbox}[breakable,colback=teal!4,colframe=teal!35!black,title={Before you move on}]
What does palpebra, palpebrae mean? (\textquotedblleft{}Eyelid\textquotedblright{}.) Say it once more.
\end{tcolorbox}
